\documentclass[11pt]{article}

\usepackage[margin=1in]{geometry}
\usepackage{amsmath,amssymb,amsthm}
\usepackage{graphicx}
\usepackage[hidelinks]{hyperref}
\usepackage{enumitem}

\newtheorem{theorem}{Theorem}
\newtheorem{lemma}{Lemma}

\newcommand{\N}{\mathbb N}
\newcommand{\Fin}{\mathrm{Fin}}
\newcommand{\I}{\mathcal I}
\newcommand{\eps}{\varepsilon}

\title{The Canonical Tensor Map for \(c_{0,\mathcal I}\)-Null Sequences}
\author{Run \texttt{fa\_banach\_001}, lane 15}
\date{2026-06-29}

\begin{document}
\maketitle

\section*{Status}

Candidate full solution, likely valid.  This packet answers Question 5.12 of
Rincon-Villamizar and Uzcategui Aylwin \cite{rvu2023}, interpreted as asking
when the canonical isometry constructed in Theorem 5.10 is onto.

\section*{Source Question}

The source paper proves that, for every Banach space \(X\) and every ideal
\(\mathcal I\) on \(\N\), the canonical map
\[
  \Phi_{\mathcal I,X}:c_{0,\mathcal I}\hat\otimes_\eps X
  \longrightarrow c_{0,\mathcal I}(X),
  \qquad
  \sum_{j=1}^m a^j\otimes x_j
  \mapsto
  \left(\sum_{j=1}^m a^j_n x_j\right)_{n\in\N}
\]
extends to an isometry into \(c_{0,\mathcal I}(X)\).  The authors then ask:

\begin{center}
  \includegraphics[width=\textwidth]{figures/open_problem_crop.png}
\end{center}

Question 5.12 asks:
\[
  \text{For which ideals }\mathcal I
  \text{ is } c_{0,\mathcal I}\hat\otimes_\eps X
  \text{ isometric to } c_{0,\mathcal I}(X)?
\]
In the context of Theorem 5.10 and Remark 5.11, this packet answers when the
source's canonical isometry is onto.

\section*{Proof Intuition}

The tensor product only sees uniform limits of finite tensor sums.  A finite
tensor sum gives a sequence whose values all lie in one finite-dimensional
subspace of \(X\), and because the sequence is bounded, its range is relatively
compact.  Uniform limits preserve relative compactness of the range.  Thus the
canonical tensor image can contain only \(c_{0,\mathcal I}(X)\)-sequences with
relatively compact range.

For \(\mathcal I=\Fin\), every norm-null sequence has relatively compact range,
so the usual \(c_0(X)\) identity gives onto-ness.  If \(X\) is
finite-dimensional, every bounded range is relatively compact, and a basis
expansion gives onto-ness for every ideal.  These are the only possibilities:
if \(\mathcal I\ne\Fin\), there is an infinite set \(A\in\mathcal I\); if \(X\)
is infinite-dimensional, Riesz's lemma gives a bounded separated sequence in
\(X\).  Placing that sequence on \(A\) gives an \(\mathcal I\)-null sequence
whose range is not relatively compact, so it cannot come from the tensor
product.

\section*{A Range Lemma}

\begin{lemma}
Every element in the range of
\[
  \Phi_{\mathcal I,X}:c_{0,\mathcal I}\hat\otimes_\eps X
  \to c_{0,\mathcal I}(X)
\]
has relatively compact range in \(X\).
\end{lemma}

\begin{proof}
For an algebraic tensor
\[
  u=\sum_{j=1}^m a^j\otimes x_j,
  \qquad a^j\in c_{0,\mathcal I},\ x_j\in X,
\]
the sequence \(\Phi_{\mathcal I,X}(u)\) takes values in the finite-dimensional
space \(E=\operatorname{span}\{x_1,\ldots,x_m\}\).  Since the scalar sequences
\(a^j\) are bounded, this range is bounded in \(E\), hence relatively compact.

Now let \(y\) be in the completed tensor image.  For every \(\eta>0\), choose
an algebraic tensor \(u\) such that
\(\|y-\Phi_{\mathcal I,X}(u)\|_\infty<\eta\).  The range of
\(\Phi_{\mathcal I,X}(u)\) is totally bounded, so finitely many \(\eta\)-balls
cover it.  The corresponding \(2\eta\)-balls cover the range of \(y\).  Thus
the range of \(y\) is totally bounded, and hence relatively compact in the
Banach space \(X\).
\end{proof}

\section*{Full Classification}

\begin{theorem}
Let \(\mathcal I\) be an ideal on \(\N\) containing \(\Fin\), and let \(X\) be
a Banach space.  The canonical isometry
\[
  \Phi_{\mathcal I,X}:c_{0,\mathcal I}\hat\otimes_\eps X
  \to c_{0,\mathcal I}(X)
\]
is onto if and only if either \(\mathcal I=\Fin\) or \(X\) is
finite-dimensional.
\end{theorem}

\begin{proof}
First suppose \(\mathcal I=\Fin\).  Then \(c_{0,\mathcal I}(X)=c_0(X)\), and
the canonical map \(c_0\hat\otimes_\eps X\to c_0(X)\) is onto.  Indeed, for
\(y=(y_n)\in c_0(X)\), the finite-support truncations
\(\sum_{n=1}^N \chi_{\{n\}}\otimes y_n\) converge to \(y\) in the sup norm.

Next suppose \(X\) is finite-dimensional.  Fix a basis \(e_1,\ldots,e_d\) of
\(X\) with coordinate functionals \(e_1^*,\ldots,e_d^*\).  If
\(y=(y_n)\in c_{0,\mathcal I}(X)\), set \(a^k_n=e_k^*(y_n)\).  Each
\(a^k\) belongs to \(c_{0,\mathcal I}\), because
\[
  \{n: |a^k_n|\geq\delta\}
  \subseteq
  \{n:\|y_n\|\geq \delta/\|e_k^*\|\}\in\mathcal I
\]
for every \(\delta>0\).  Hence
\[
  y=\Phi_{\mathcal I,X}\left(\sum_{k=1}^d a^k\otimes e_k\right),
\]
so \(\Phi_{\mathcal I,X}\) is onto.

Conversely, assume that \(\mathcal I\ne\Fin\) and \(X\) is infinite-dimensional.
Choose an infinite set \(A=\{n_1<n_2<\cdots\}\in\mathcal I\).  By Riesz's
lemma, there is a bounded sequence \((x_k)\) in \(X\) with no relatively compact
subsequence; for instance one may take a sequence in the unit sphere separated
by a fixed positive constant.

Define \(y=(y_n)\in\ell_\infty(X)\) by
\[
  y_{n_k}=x_k\quad(k\geq1),
  \qquad
  y_n=0\quad(n\notin A).
\]
For every \(\delta>0\), the set \(\{n:\|y_n\|\geq\delta\}\) is a subset of
\(A\), hence belongs to \(\mathcal I\).  Thus \(y\in c_{0,\mathcal I}(X)\).
However, the range of \(y\) contains the non-relatively-compact set
\(\{x_k:k\geq1\}\).  By the range lemma, \(y\) cannot lie in the range of
\(\Phi_{\mathcal I,X}\).  Therefore \(\Phi_{\mathcal I,X}\) is not onto.
\end{proof}

\section*{Verification Notes}

No computational verification is involved.  The proof uses only:
\begin{itemize}
\item finite-dimensional compactness of bounded sets;
\item the finite-support truncation proof of
\(c_0\hat\otimes_\eps X\cong c_0(X)\);
\item Riesz's lemma to produce a bounded non-totally-bounded sequence in any
infinite-dimensional Banach space;
\item the source paper's canonical embedding from Theorem 5.10.
\end{itemize}

The statement is for the canonical map \(\Phi_{\mathcal I,X}\).  It does not
claim that no noncanonical isometric isomorphism can exist between the two
Banach spaces in exceptional cases.

\section*{Novelty Check}

Local run-index searches for \(\texttt{2309.08076}\),
\(\texttt{Question 5.12}\), \(\texttt{injective tensor}\),
\(\texttt{c\_\{0,I\}(X)}\), \(\texttt{relatively compact range}\), and
\(\texttt{finite-dimensional}\) found no previous full classification packet.
The earlier lane-15 packet
\(\texttt{2309.08076\_nonfin\_tensor\_obstruction\_no\_complemented\_c0}\)
is superseded by this full statement because it missed the finite-dimensional
positive case.

Bounded web searches on 2026-06-29 for exact and close phrases including
\[
\texttt{"c\_\{0,}\backslash\texttt{mathcal I\}" "relatively compact range"
"injective tensor"},
\]
\[
\texttt{"c\_\{0,I\}(X)" "injective tensor" "finite-dimensional"},
\quad
\texttt{"Question 5.12" "c\_\{0,I\}" "c\_\{0,I\}(X)"},
\]
and the source title together with \(\texttt{Question 5.12}\) found no explicit
later answer in this pass.

\section*{Limitations}

This answers the source question for the canonical isometry constructed in the
paper.  If one reads Question 5.12 as asking about arbitrary isometric
isomorphism of the two Banach spaces, rather than onto-ness of the canonical
map, that broader formulation is not claimed here.

\section*{Human Review Recommendation}

Likely valid.  Review should focus on the intended reading of Question 5.12
and on the range lemma.

\begin{thebibliography}{9}

\bibitem{rvu2023}
M. A. Rincon-Villamizar and C. Uzcategui Aylwin,
\emph{Banach spaces of \(\mathcal I\)-convergent sequences},
arXiv:2309.08076, 2023.

\end{thebibliography}

\end{document}
